We trained the pretrained YOLOv5 model using a learning rate of 0.01 for a total of 100 epochs. Table \ref{tab:detection-result} provides a summary of the resulting performance metrics, including the average mean precision at 0.5 IoU threshold (mAP@0.5), average mean precision at IoU threshold between 0.5 and 0.95 (mAP@[0.5:0.95]), best precision, and best recall. Following the overall results, we proceed to analyze the outcomes of each group of augmentations, namely Geometric Augmentation, Color-Space Augmentation, Cutout, and RandAugment.

\subsubsection{Overall}
To summarize, most of the geometric augmentations negatively impacted the model's performance, with the exception of CenterCrop. CenterCrop achieved the highest mAP@0.5 and mAP@0.5:0.95, respectively with values of 0.3436 and 0.5214, which is nearly two times better than the baseline performance. On the other hand, color-space augmentations generally improved the model's performance, with the best augmentation technique being ToSepia, which scores 24\%  in mAP@[0.5:0.95], a notable 6.29\% improvement over the baseline. Cutout and GridDropout had a notable negative impact on the model's performance, resulting in a decrease of 2\% and 7.5\%, respectively. However, the worst-performing augmentation technique was RandAugment, which produced a staggering 10\% decrease in the mAP.

\begin{table}[!htb]
    \centering
    \caption{Results of augmentations on YOLOv5s at fine-tuned learning rate $10^{-3}$}    \label{tab:detection-result}
    \begin{tabular}{|l|llll|}
    \hline
        \textbf{Method} & \textbf{\makecell{mAP@[0.5:0.95]}} & \textbf{\makecell{mAP@0.5}} & \textbf{\makecell{Precision}} & \textbf{\makecell{Recall}} \\ \hline
        Baseline & 0.1773 & 0.2879 & 0.9298 & 0.6391\\ 
        \hline
        CenterCrop 0.2 & \textbf{0.3436} & \textbf{0.5214} & 0.7719 & 0.6865 \\ 
        Flipping & 0.1584 & 0.2772 & \textbf{0.9576} & \textbf{0.6931} \\ 
        Motion Blur & 0.1517 & 0.2256 & 0.4693 & 0.2177 \\ 
        Perspective & 0.1244 & 0.1999 & 0.2695 & 0.2256 \\
        Rotation & 0.1427 & 0.2303 & 0.3325 & 0.2665 \\ 
        \hline
        CLAHE & 0.2152 & 0.3369 & 0.7346 & 0.4768 \\ 
        ChannelShuffle & 0.235 & 0.3619 & 0.7651 & 0.575 \\ 
        GaussNoise & 0.2077 & 0.3236 & 0.7858 & 0.525 \\ 
        Solarize & 0.1612 & 0.2683 & 0.3872 & 0.273 \\
        Sharpen & 0.2402 & 0.3604 & 0.7413 & 0.4545 \\ 
        ToGray & 0.2394 & \textbf{0.363} & \textbf{0.7741} & \textbf{0.5432} \\ 
        ToSepia & \textbf{0.246} & 0.3757 & 0.7404 & 0.5437 \\ 
        \hline
        Cutout & \textbf{0.1503} & \textbf{0.2481} & 0.4137 & 0.2427 \\ 
        GridDropout & 0.1023 & 0.1535 & 0.3643 & 0.1549 \\ 
        \hline
        RandAugment & 0.07167 & 0.1347& 0.5115& 0.1305 \\
    \hline
    \end{tabular}
\end{table}

\subsubsection{Geometric Augmentation}\label{result:dec:geo}
The mAP@[0.5:0.95] metric for the geometric augmentation technique over 100 epochs is visually represented in Figure \ref{fig:yolo_geo}. As demonstrated in the results, CenterCrop augmentation technique outperforms all other augmentation techniques. This observation confirms our initial assumption that an augmentation technique that ensures center invariance can considerably enhance the model's performance in terms of both bounding box detection and class prediction, as depicted in Figure \ref{fig:geo_loss}.

\begin{figure}[!h]
	 \centerline{\includegraphics[width=0.9\textwidth]{src/Results/detection/yolo_geo.png}}
	 {\caption{Mean average precision in threshold interval between 0.5 and 0.95 of YOLOv5s with geometric augmentations}\label{fig:yolo_geo}}
\end{figure}

Although it is reasonable to assume that motion blur would not improve a model's performance due to the added noise, the fact that the Rotation, Flipping, and Perspective augmentations decreased the overall performance of the model is counter-intuitive. These transformations are generally expected to increase the robustness of the training set, but they may actually have an adverse effect on the model's ability to learn.

One possible explanation for this unexpected outcome is that YOLOv5 is a bounding box detection model. A modification in the geometric shape of the object can result in a more varied squared shape of the bounding box, which can slow down the training convergence. This assumption is supported by the results shown in figure \ref{fig:geo_loss}, where we observe that although the classification loss did not substantially increase with the augmentations (except for Perspective), the box loss of Flipping, Perspective, and Rotation decreased much slower than the baseline. Hence, the architecture of the YOLOv5 model needs to be considered while selecting augmentation techniques for improving the model's performance.

\begin{figure}
\hfill
\subfigure[Box loss]{\includegraphics[width=7cm]{src/Results/detection/geo_box_loss.png}}
\hfill
\subfigure[Classification loss]{\includegraphics[width=7cm]{src/Results/detection/geo_class_loss.png}}
\hfill
\caption{Illustration of (a) box loss and (b) classification loss in geometric augmentation group}
\label{fig:geo_loss}
\end{figure}

\subsubsection{Color-Space Augmentation}

\begin{figure}[!h]
	 \centerline{\includegraphics[width=0.9\textwidth]{src/Results/detection/yolo_color.png}}
	 {\caption{Mean average precision of YOLOv5s with Color-Space Augmentations}\label{fig:yolo_color}}
\end{figure}

Figure \ref{fig:yolo_color} provides a visualization of the mAP@[0.5:0.95] metric for color-space augmentations. In general, all color-space augmentation techniques significantly improve the performance of YOLOv5.

The plot in Figure \ref{fig:color_loss} illustrates the trade-off between classification loss and box-loss in color-space augmentation. The extra sharpness and contrast introduced by this augmentation can significantly improve the model's ability to predict the correct class label. However, this comes at a cost of decreased accuracy in detecting the bounding box. This indicates that the model with color-augmentation could also learn features of the background to aid in object classification.

\begin{figure}[!h]
\hfill
\subfigure[Box loss]{\includegraphics[width=7cm]{src/Results/detection/color_box_loss.png}}
\hfill
\subfigure[Classification loss]{\includegraphics[width=7cm]{src/Results/detection/color_class_loss.png}}
\hfill
\caption{Illustration of (a) box loss and (b) classification loss in Color-Space Augmentations}
\label{fig:color_loss}
\end{figure}

While CLAHE and Sharpen work by enhancing the contrast and brightness of images, making it easier for the model to detect and classify objects, the fact that ToSepia and ToGray augmentations also improve performance is noteworthy. Interestingly, ToSepia outperforms all other color-space augmentations, indicating that reducing color information can enhance the model's detection ability. By reducing the variances in the three channels and consolidating them into one dimension, the model can focus more on the object's shape rather than color, resulting in improved performance.

\subsubsection{Cutout and RandAugment}
Figure \ref{fig:yolo_cutout} shows that the performance of the model is reduced by both Cutout and GridDropout. Furthermore, RandAugment causes the most significant reduction in performance, with the largest margin observed among all augmentation techniques.

\begin{figure}[!h]
	 \centerline{\includegraphics[width=0.9\textwidth]{src/Results/detection/yolo_cutout.png}}
	 {\caption{Mean average precision of YOLOv5s with Cutout Augmentations and RandAugment}\label{fig:yolo_cutout}}
\end{figure}

The reasons for the decrease in performance caused by Cutout and Dropout can be relatively intuitive, as they are non-label-preserving augmentation techniques that may remove crucial information required for the model to detect and classify objects accurately. Figure \ref{fig:cutout_loss} demonstrates that Cutout produces a significant increase in the box loss, even more so than GridDropout, which actually removes more information. Meanwhile, RandAugment performs poorly in both classification and box losses, introducing too many inconsistencies into the model and hindering its ability to identify crucial features.

\begin{figure}[!h]
\hfill
\subfigure[Box loss]{\includegraphics[width=7cm]{src/Results/detection/cutout_box_loss.png}}
\hfill
\subfigure[Classification loss]{\includegraphics[width=7cm]{src/Results/detection/cutout_class_loss.png}}
\hfill
\caption{Illustration of (a) box loss and (b) classification loss in Cutout augmentations and RandAugment}
\label{fig:cutout_loss}
\end{figure}